%%%%%%%%%%%%%%%%%%%%%%%%%%%%%%%%%%%%%%%%%%%%%%%%%%%%%%%%%%%%%%%%%%%%%%%%%%%%%%%%
%2345678901234567890123456789012345678901234567890123456789012345678901234567890
%        1         2         3         4         5         6         7         8

\documentclass[a4paper, 10pt, conference]{IEEEtran} % Use this line for a4 paper
\IEEEoverridecommandlockouts % This command is only needed if 

\usepackage{graphics} % for pdf, bitmapped graphics files
\usepackage{epsfig} % for postscript graphics files
\usepackage{mathptmx} % assumes new font selection scheme installed
\usepackage{times} % assumes new font selection scheme installed
\usepackage{amsmath} % assumes amsmath package installed
\usepackage{amssymb}  % assumes amsmath package installed

% Custom packages
\usepackage{algorithm}
\usepackage{algpseudocode} % pour avoir l'environnement algorithmic

\title{\LARGE \bf
--- Electronical Engineering Note --- \\
Engineering of a 94kHz 15A sensorless quadricopter ESC
}

\author{Stepane Binon and Guilhem Andrieu% <-this % stops a space
	\thanks{Bare Metal Foundry}
	\thanks{\tt\small stepane.binon@gmail.com}
}

\begin{document}

\maketitle
\thispagestyle{empty}
\pagestyle{empty}

%%%%%%%%%%%%%%%%%%%%%%%%%%%%%%%%%%%%%%%%%%%%%%%%%%%%%%%%%%%%%%%%%%%%%%%%%%%%%%%%
\begin{abstract}
To come
\end{abstract}


%%%%%%%%%%%%%%%%%%%%%%%%%%%%%%%%%%%%%%%%%%%%%%%%%%%%%%%%%%%%%%%%%%%%%%%%%%%%%%%%
\section{Introduction}

The goal of this document is to sumarize the though process, the choices and calculations made during the design of an ESC aimed at being used in an autonomous race quadricopter. 

This project is part of the Bare Metal Foundry, an association which aims at having serious fun with engineering challenges and is also our first attemps at PCB design, both as engineers and association. In this context, the Kururugi-ESC has a central place as a vector of learing and as one of the most important part of a quadricopter, the ESC being the deepest loop of control, the fastest, and thus the deadliest in case of crash. 

We wanted to do this design from scratch to, in addition of learning, gain complete control of the hardware (h/w) and matching the performances we needed. Hence, for now, we do not aim to make it compatible with any ESC software existing as we will develop a custom one in C++, this will be the subject of an article to come. For Version 1.0 of this project, it is aimed to contribute in the following points:
\begin{itemize}
    \item The production of a 24 to 94kHz, 15A continuous, 95\%D (Duty Cycle), 4S to 6S, and sensorless ESC that can work both with Six-Step Control (SSC) and Field-Oriented Control (FOC);
    \item The presentation of measure and results obtained;
    \item A complete guideline on how to design a such an ESC PCB, that we hope is state of the art.
\end{itemize}

The Version 1.0 of the Kururugi-ESC will be oriented toward getting it working whereas next versions goals are aimed toward research. One particular point we are interested in is to have the SSC working on a 128kHz PWM, which would bring challenges both for the electronical and coding aspects.

The rest of this paper is organized as follows. Section II will be about the schmatics and selection of components. Section III will present the layout. Section II will exposes results


%%%%%%%%%%%%%%%%%%%%%%%%%%%%%%%%%%%%%%%%%%%%%%%%%%%%%%%%%%%%%%%%%%%%%%%%%%%%%%%%
\section{Schematic and choice of components}

An ESC is composed of N main sectionsparts
\begin{itemize}
    \item An MCU (Micro Controller Unit), responsible for the overall control of the behaviour, mainly communication with the Flight Controller (FC) and control of the power section behaviour;
    \item Analog to Digital Converters (ADC) and Operational Amplificator (OpAmp), which are used to translate analog signal as tension and current on each phase to the digital language of the MCU. As space is constrained on this PCB, the MCU internal ADC and OpAmp will be used;
    \item The power section, composed of the 3 half bridge needed to control the Brushless Direct Current Motor (BLDC), get the necessary analog outputs for the MCU and also power the ICs present on the PCB. This is from far the most complexe part of an ADC as it require fine control over very important electron flows;
\end{itemize}

\subsection{Half bridge}

To control the BLDC motor by mimicing electronically the behaviour of an MCC motor, an ESC relies on 3 phases that exictes the different phases at precise timings. This paper does not sumarize the theory behind it. The first role of an ESC is to ensure proper commutation to keep the motor running. Its second role it to control at which speed the motor rotate. 

The design of the 3 half bridge needed is the most critical part as Gate-Drive Domain and Logic Domain collide with the Power Domain. Among other, these are the main challenges:

\begin{itemize}
    \item Ensuring proper conduction of the battery power to the motor cables.
    \item Ensuring minimal EMI creation during the high-speed switching of mostfet.
    \item Ensuring proper sensing of the voltage and current on the motor phases, this is critical for sensorless FOC. The low power sensing must not be flowed by high power EMI noise. Ensuring this sensing, which mainly relies on resisor, does not creates too much heat.
    \item Ensuring that the Gate capacitance is filled fast enough for high speed switching and low switching losses
    \item Ensuring that their is no creep and false trigger in the Gate capacitance
    \item Preventing shoot through caused by simultaneous trigger of MOSFET
    \item Ensuring all components are cooled enough and that the ESC overall temperature stay tolerable
    \item Ensuring that the boost capacitor of the gate driver is filled fast enough
    \item Ensuring that the gate driver capacitance is enough for all workload and temperatures
\end{itemize}

To chose a mosfet, over that maximum current admissible at maximal running temperature, two parameter are of the outmost importance:
\begin{itemize}
    \item $Q_g$ which defines the gate capacitance. An increase of gate capacitance means that more power will be need to both turn on and off the mosfet, wich creates more heat and longer switching time.
    \item $R_{DS,on}$ which defines the resistance of the mosfet when on. This is used to compute the heat created by the mosfet and is important at high duty cycles. 
\end{itemize}
At low PWM frequency, conduction losses are more important than swithing losses, this tendency inverted at high frequency PWM and much more heat is create overall.

Conduction losses per mosfet can be calculated using:
\[
P_{cond} = R_{DS,on} * I_{max} * D_{max}
\]

Swithing losses per mosfet can be calculated using:
\[
P_{sw} = U_{max} * I_{max} * (t_{rise} + t_{fall}) * F_{ws,max}
\]

The selection of the MOSFET will define all the Half-Bridge. For the Kururugi-ESC-V1.0, six Texas Instruments CSD18540Q5B will be used. 





\end{document}